\chapter{Conclusion and future work}\label{ch:conclusion}

\section{Conclusion}\label{sec:summary}

\fcsg{} splits an image into three frequency bands and routes each band, per
pixel, to two of three shared convolutional experts. With 97,088 parameters
and 9.39 GFLOPs, it passed every check before training (O1). Our FFDNet code
matched the published CBSD68 result to within 0.01\dB{} (O2). Under one
matched composite-degradation protocol, \fcsg{} reached 27.24\dB{}. That is
0.59\dB{} below DnCNN and 0.95\dB{} below blind FFDNet, with 82.6\% and
88.6\% fewer parameters. It was also the slowest model on a Tesla T4 (O3).
Its gate varies across the image and follows the noise level, with
correlations up to 0.81. The trained model loses 0.37\dB{} when we override
the gate's choices (O4). Frequency-band routing gave a small model behaviour
we can inspect. It did not beat simple baselines on accuracy or speed.

\section{Limitations}\label{sec:limitations}

We trained each model once, and all results come from the validation split,
with no untouched test set. The degradation is synthetic. We did not test on
real archival photographs or run a human study. The routing analysis rests on
ten images and exploratory correlations. The interventions measure how much
the trained model leans on its gate, not what a gate adds during training. We
measured latency on one GPU without an operator profile. Every expert always
runs, so sparse routing saves no computation.

\section{Future work}\label{sec:future}

The first experiment to run is a matched training ablation. It would compare
learned routing with fixed uniform weights, from the same initialisation and
with several seeds. Further ablations would test the kernel sizes, a
per-image gate, plain concatenation in place of SE fusion, and hard masks.
The model relies mainly on which experts it picks, so skipping unselected
experts could save up to a third of the expert cost. An operator profile
should come first, along with latency measurements on CPU and mobile
hardware. The parameter budget is fifty times what the model uses, so
widening the experts could close the capacity gap, provided cheaper expert
evaluation frees up compute. Last, any claim about archival restoration needs
real archival images, a controlled degradation sweep and a held-out test set.
